\chapter{Water-soluble vitamins: C and the B group}

\section{Learning objectives}
Explain the biochemical roles of vitamin C and the B vitamins; identify high-risk dietary and medical contexts; recognize the neurological importance of B12; and differentiate physiological replacement from pharmacological dosing.

\section{Vitamin C}
Vitamin C is required for collagen synthesis, wound healing, antioxidant protection, and improved absorption of non-heme iron from plant foods. Citrus, guava, berries, peppers, tomatoes, broccoli, and potatoes contribute vitamin C. Deficiency causes scurvy, with bleeding gums, poor wound healing, bruising, and fatigue. Large supplemental doses can cause gastrointestinal upset and may increase kidney-stone risk in people with pre-existing hyperoxaluria or renal disorders. Vitamin C does not reliably prevent the common cold, although regular use may modestly shorten symptoms for some people. \cite{odsC}

Ascorbate is a reducing agent and cofactor for enzymes involved in collagen maturation, carnitine synthesis, and catecholamine metabolism. In U.S. dietary reference values, smokers have a higher recommended intake; severe dietary restriction and malabsorption can increase deficiency risk. Low vitamin C status can occur in people receiving dialysis, but supplementation should be clinician directed because excess ascorbate can increase oxalate. Scurvy can present with perifollicular hemorrhage, corkscrew hairs, gingival disease, arthralgia, anemia, and impaired wound healing. \cite{odsC}

\section{The B vitamins}
The following summary should be read with the corresponding NIH health-professional fact sheets, because dose, deficiency risk, and toxicity vary by age, pregnancy, disease, and formulation. \cite{odsList}
\begin{longtable}{p{0.18\textwidth}p{0.36\textwidth}p{0.34\textwidth}}
\caption{B-vitamin functions and food patterns}\\
\toprule
Nutrient & Main roles & Useful sources or cautions \\
\midrule
\endfirsthead
\toprule Nutrient & Main roles & Useful sources or cautions \\
\midrule
\endhead
B1 (thiamin) & Carbohydrate metabolism and nerve function. & Whole grains, legumes, pork, nuts; risk rises with severe alcohol use disorder. \cite{odsThiamin}\\
B2 (riboflavin) & Energy metabolism and cellular oxidation-reduction. & Dairy, eggs, meat, almonds, fortified grains. \cite{odsRiboflavin}\\
B3 (niacin) & Energy metabolism, skin, nerves, and digestion. & Meat, fish, peanuts, legumes, fortified grains; pharmacologic doses can flush or injure the liver. \cite{odsNiacin}\\
B5 (pantothenic acid) & Coenzyme A and fatty-acid metabolism. & Widely distributed; deficiency is rare. \cite{odsPantothenic}\\
B6 (pyridoxine family) & Amino-acid metabolism, neurotransmitters, and hemoglobin. & Poultry, fish, potatoes, bananas, chickpeas; chronic high doses can cause neuropathy. \cite{odsB6}\\
B7 (biotin) & Carboxylation reactions in metabolism. & Eggs, nuts, legumes; high-dose biotin can distort thyroid, cardiac, and other laboratory tests. \cite{odsBiotin}\\
B9 (folate/folic acid) & DNA synthesis and cell division; red-cell formation. & Leafy greens, legumes, citrus, fortified grains. Folic acid before conception and during early pregnancy reduces neural-tube-defect risk; large amounts can correct B12-related anemia without treating neurological damage. \cite{odsFolate}\\
B12 (cobalamin) & DNA synthesis, red-cell formation, and nerve maintenance. & Animal foods or fortified foods; absorption depends on stomach and intrinsic factor. Vegans need a reliable source, such as fortified foods or a supplement. \cite{odsB12}\\
\bottomrule
\end{longtable}

Thiamin deficiency impairs oxidative metabolism and can cause beriberi or Wernicke--Korsakoff syndrome. In a person at high risk of deficiency, thiamin should be given promptly when clinically indicated, but emergency glucose for hypoglycemia should not be delayed while waiting for thiamin. Niacin deficiency produces pellagra classically characterized by dermatitis, diarrhea, and dementia. B6 deficiency can cause anemia or neuropathy, but high supplemental exposure can itself cause sensory neuropathy. \cite{odsThiamin,thiaminGlucose,odsNiacin,odsB6}

\subsection{Clinical profiles of the B vitamins}
\begin{description}
\item[Thiamin.] Consider in severe alcohol use disorder, prolonged vomiting, starvation, bariatric surgery, and refeeding. Neurologic or cardiac features require urgent assessment; normal routine blood tests do not exclude Wernicke encephalopathy. \cite{odsThiamin}
\item[Riboflavin.] Deficiency may cause angular stomatitis, cheilosis, glossitis, seborrheic dermatitis, and anemia, usually with other nutritional deficiencies. Assessment is clinical and dietary; isolated severe deficiency is uncommon. \cite{odsRiboflavin}
\item[Niacin.] Pellagra risk increases with malabsorption, alcoholism, carcinoid syndrome, Hartnup disease, and isoniazid exposure. Pharmacological nicotinic acid may be prescribed for dyslipidemia, but added cardiovascular benefit with statin therapy has not been established; high doses require clinician oversight and monitoring. \cite{odsNiacin}
\item[Pantothenic acid.] Deficiency is rare because the nutrient is widely distributed. Burning feet and fatigue are nonspecific and should not be attributed to pantothenate without a compatible clinical context. \cite{odsPantothenic}
\item[Vitamin B6.] Pyridoxal phosphate supports transamination, neurotransmitter synthesis, heme production, and glycogen metabolism. Both deficiency and chronic supplemental excess can cause neuropathy. \cite{odsB6}
\item[Biotin.] Deficiency is uncommon. High-dose biotin may interfere with immunoassays; patients should disclose use before laboratory testing. \cite{odsBiotin}
\item[Folate.] Deficiency produces impaired DNA synthesis and megaloblastic anemia. Periconceptional folic acid prevents neural-tube defects, while unexplained macrocytosis should prompt consideration of B12 deficiency before folate-only treatment. \cite{odsFolate}
\item[B12.] Deficiency may cause macrocytosis, glossitis, neuropathy, ataxia, or cognitive symptoms, sometimes without anemia. Evaluate intake, gastric and ileal disease, intrinsic factor, medicines, and kidney function. \cite{odsB12}
\end{description}

Folate is central to nucleotide synthesis and methyl-group transfer; deficiency produces megaloblastic anemia and, in pregnancy, increases neural-tube-defect risk. B12 is required for methionine synthase and methylmalonyl-CoA mutase. Serum B12 can be misleading; methylmalonic acid and homocysteine may help in selected cases, interpreted with renal function and clinical context. \cite{odsB12,odsFolate}

\section{Folate and B12 deserve special attention}
Folate needs rise in pregnancy, and supplementation should follow local prenatal guidance. B12 deficiency may cause anemia, numbness, balance problems, or cognitive changes, and neurological injury can occur even without anemia. Risk is higher with vegan diets, pernicious anemia, gastric or intestinal surgery, and long-term use of some acid-suppressing medicines or metformin. Diagnosis should not be based on symptoms alone. \cite{odsFolate,odsB12}

\warning{Do not use high-dose B-complex products as a general energy treatment. Correcting a deficiency may help; extra B vitamins do not create energy in the absence of inadequate intake.}

\section{Clinical approach}
Macrocytosis may reflect folate or B12 deficiency, alcohol, liver disease, hypothyroidism, reticulocytosis, marrow disorders, or medications. Neuropathy may reflect B12 deficiency, diabetes, alcohol, toxins, compression, or other neurological disease. High folic-acid intake can correct anemia associated with B12 deficiency without treating its neurological effects; whether it masks deficiency enough to worsen neurological outcomes remains uncertain. Consider B12 risk before treating unexplained macrocytosis with folate alone. High-dose biotin can interfere with laboratory assays, including some thyroid and cardiac tests. \cite{odsFolate,odsB12,odsBiotin}

\section{Key points}
Water solubility does not eliminate toxicity; dose and duration determine risk. Replace documented deficiencies while investigating the dietary, gastrointestinal, medication, or systemic cause, and arrange follow-up to confirm response.
